% =========================================================================
% Memoria: adaptación y evaluación en HotpotQA del sistema de recuperación
% multisalto sobre grafo de aserciones reificadas con memoria canónica.
% Compilar: tectonic metodologia_hotpotqa.tex   (o pdflatex, dos pasadas)
% Complementa a metodologia_2wiki.tex (método completo y caso clínico).
% =========================================================================
\documentclass[11pt,a4paper]{article}
\usepackage[T1]{fontenc}
\usepackage[utf8]{inputenc}
\usepackage[spanish,es-noshorthands,es-tabla]{babel}
\usepackage{lmodern}
\usepackage[margin=2.4cm]{geometry}
\usepackage{amsmath,amssymb,mathtools}
\usepackage{booktabs}
\usepackage{array}
\usepackage{caption}
\usepackage{enumitem}
\usepackage[hidelinks]{hyperref}

\newcommand{\sys}{\textsc{CatRAG-Reif}}
\newcommand{\hippo}{HippoRAG~2}
\newcommand{\gl}{\guillemotleft}
\newcommand{\gr}{\guillemotright}
\DeclareMathOperator*{\argmax}{arg\,max}
\setlength{\parskip}{2pt}

\title{Generalización sin ajuste de la recuperación multisalto sobre un grafo
de aserciones reificadas con memoria canónica de entidades:\\
adaptación del método y evaluación en HotpotQA}
\author{}
\date{30 de agosto de 2026}

\begin{document}
\maketitle
\vspace{-2.5em}

\begin{abstract}
\noindent
Un sistema de recuperación multisalto diseñado y calibrado sobre
2WikiMultihopQA ---aserciones reificadas como nodos, memoria canónica de
entidades derivada del corpus, paseo con reinicio personalizado y un analista
de la pregunta que planifica la cadena--- se traslada a HotpotQA con la
configuración \emph{congelada}: ni un dial, ni un umbral, ni una demostración
cambiados. Para que el traslado sea honesto se retiran antes del flujo de
consulta las tres reglas específicas de 2Wiki (enrutador léxico, salto de la
selección de conjunto en comparación, recorte de menciones por mayúsculas
como única vía de enlace), se cura la memoria canónica con decisiones
registradas y auditables, y se generaliza el banco de pruebas a un esquema
común. Sobre el subconjunto de reproducción de HotpotQA usado por \hippo{} y
CatRAG ($1\,000$ preguntas de nivel \emph{hard}, corpus de $9\,811$ pasajes) y
con los mismos modelos que esos trabajos (\texttt{gpt-4o-mini},
\texttt{text-embedding-3-small}), la medición única del banco da
$\mathrm{R@}5$ $94{,}7$ y cadena completa a $k{=}5$ ($\mathrm{FC@}5$)
$90{,}2$, frente a $89{,}5$ / $80{,}4$ de CatRAG y $87{,}1$ / $75{,}5$ de
\hippo{} publicados con idéntico equipamiento; pareado frente a BM25 sobre el
mismo corpus, $+40{,}7$ puntos de $\mathrm{FC@}5$ (IC del $95\%$
$[+37{,}1,+44{,}1]$). En el sondeo retenido, las ablaciones atribuyen la
ganancia a las mismas piezas que en 2Wiki: el salto dirigido ($+2{,}0$ de
$\mathrm{FC@}5$, significativo) y la selección de conjunto ($+4{,}0$ de
$\mathrm{R@}2$). El análisis de fallos ---$98$ de $1\,000$, $96$ de puente, en
$90$ de ellos falta un solo pasaje--- no muestra un patrón nuevo respecto a
2Wiki. Se documentan el protocolo, los costes ($\approx 5$ USD en total) y
las limitaciones: sin identificadores externos para auditar la memoria, sin
reproducción propia de \hippo{} sobre este corpus, y una sola familia de
modelos.
\end{abstract}

% =========================================================================
\section{Introducción}
\label{sec:intro}

La pregunta que responde esta memoria es estrecha y verificable: un método
de recuperación multisalto que rinde bien en el conjunto sobre el que se
diseñó, ¿rinde también en otro conjunto \emph{sin tocarlo}? La tentación
habitual ---reajustar diales, reescribir consignas, añadir reglas para el
conjunto nuevo--- convierte cada resultado en una anécdota. Aquí se hace lo
contrario: se fija la configuración, se escribe de antemano qué se espera
(hipótesis H1 del documento de método) y se mide una sola vez.

El sistema, descrito en detalle en la memoria de método sobre
2WikiMultihopQA~\cite{wiki2} (en adelante, \emph{el documento de 2Wiki}),
tiene tres ideas: (i) la unidad semántica del índice es la \emph{aserción
reificada} ---un hecho promovido a nodo con su frase autocontenida y un rol
por participante--- en lugar del triplete; (ii) las entidades se resuelven
en indexación contra una \emph{memoria canónica} derivada del propio corpus
(fichas por mención, dos canales de similitud, restricciones duras, juez y
cierre transitivo), de modo que el grafo no necesita aristas de sinonimia y
cada entidad tiene un \emph{pasaje propio}; (iii) en consulta, un
\emph{analista} planifica la cadena de huecos, un paseo con reinicio
personalizado ordena los pasajes y un \emph{salto dirigido} rellena los
huecos que el primer paseo deja abiertos. Sobre 2Wiki, con
\texttt{gpt-4o-mini}, esa configuración (v4) da $\mathrm{R@}5$ $98{,}1$ y
$\mathrm{FC@}5$ $96{,}0$.

HotpotQA~\cite{hotpotqa} es el contraste natural: preguntas escritas por
personas y no generadas por plantillas, con menciones descriptivas y en
minúsculas, dos tipos (puente y comparación) en lugar de cuatro, y sin los
identificadores de Wikidata que en 2Wiki permitieron auditar la memoria
canónica. Es además el conjunto en el que \hippo{}~\cite{hipporag2} y
CatRAG~\cite{catrag} publican sus cifras con los mismos modelos que usamos,
lo que hace posible una comparación directa sin reproducir nada.

Las contribuciones de esta memoria son cuatro. Primera, el inventario
explícito de lo que en la versión evaluada en 2Wiki era específico del
conjunto y su sustitución por mecanismos generales
(sección~\ref{sec:adaptacion}). Segunda, la curación auditada de la memoria
canónica ---uniones con procedencia, juez para las uniones de título,
homónimos con desambiguador--- validada en 2Wiki con el oro externo antes de
aplicarla aquí (sección~\ref{sec:memoria}). Tercera, el protocolo de
generalización con sondeo retenido, configuración congelada y medición
única, y sus resultados (secciones~\ref{sec:protocolo}
y~\ref{sec:resultados}). Cuarta, un análisis de fallos que sitúa el residuo
y una discusión de validez que dice qué no se ha medido
(secciones~\ref{sec:fallos} y~\ref{sec:validez}).

% =========================================================================
\section{HotpotQA y el banco de pruebas}
\label{sec:datos}

\paragraph{El conjunto.}
HotpotQA~\cite{hotpotqa} contiene $113$k preguntas construidas por
anotadores a partir de \emph{pares} de párrafos introductorios de Wikipedia.
En las preguntas de \emph{puente} ($\approx 80\%$) el segundo párrafo solo
se identifica a través de una entidad que el primero menciona (\gl{}What 2003
Christmas-themed romantic comedy did Gregor Fisher have a role in?\gr{}: de
la biografía de Fisher a \emph{Love Actually}); en las de \emph{comparación}
se contrastan dos entidades nombradas (\gl{}Which rock band chose its name by
drawing it out of a hat, Switchfoot or Midnight Oil?\gr{}). Cada pregunta
trae, en el ajuste \emph{distractor}, diez párrafos ---los dos oro y ocho
distractores recuperados por TF-IDF--- y las frases de apoyo
(\texttt{supporting\_facts}) a nivel de frase; un campo \texttt{level}
(\emph{easy}, \emph{medium}, \emph{hard}) gradúa la dificultad.

\paragraph{El banco.}
Seguimos el protocolo de \hippo{} y CatRAG: $1\,000$ preguntas del
\emph{dev} (todas de nivel \emph{hard}; $811$ de puente, $189$ de
comparación) y un corpus cerrado formado por la unión de sus contextos,
$9\,811$ pasajes con título único. La recuperación es \emph{abierta} contra
los $9\,811$, no contra los diez de cada pregunta; el oro son los dos títulos
de \texttt{supporting\_facts}. El subconjunto está vendorizado con el código
de \hippo{} y es el mismo que usa CatRAG, de modo que sus cifras publicadas
con \texttt{gpt-4o-mini} + \texttt{text-embedding-3-small} son referencia
directa (tabla~\ref{tab:banco}).

\paragraph{Sondeo y calibración.}
Las decisiones ---ninguna en este trabajo, pero el protocolo las prevé--- se
toman sobre un \emph{sondeo} de $200$ preguntas del conjunto de
entrenamiento (disjunto del banco), estratificado en $100$ de puente y $100$
de comparación, restringido al nivel \emph{hard} para igualar la dificultad
del banco, con su mini-corpus ($1\,990$ pasajes: unión de sus contextos). Las
$15\,461$ preguntas \emph{hard} restantes del entrenamiento quedan como
cantera de calibración. El conjunto de entrenamiento se obtuvo del espejo de
HuggingFace en formato parquet (el servidor oficial no respondía) y se
convirtió al esquema original.

\paragraph{Métricas.}
Para cada pregunta y $k\in\{2,5,10,20\}$: $\mathrm{recall@}k$, fracción de
títulos oro entre los $k$ primeros, y $\mathrm{FC@}k$ (\emph{full-chain
retrieval}), que vale $1$ solo si \emph{todos} los títulos oro están entre los
$k$ primeros ---la métrica que exige la cadena completa y la que CatRAG
llama FCR---. Los agregados llevan intervalos de confianza del $95\%$ por
\emph{bootstrap} ($2\,000$ remuestreos) y las comparaciones entre sistemas
sobre las mismas preguntas son \emph{pareadas}: diferencia media con su
intervalo y recuento de preguntas ganadas y perdidas. Con dos oros por
pregunta, $\mathrm{FC@}2$ es alcanzable en todos los tipos (a diferencia de
la comparación-puente de 2Wiki).

\paragraph{Diferencias con 2Wiki que importan al método.}
La tabla~\ref{tab:diferencias} las resume. Tres pesan: las preguntas son más
largas y libres (16 palabras de media frente a 12), con $35$ de las $1\,000$
sin ninguna palabra en mayúscula tras la primera, lo que desarma cualquier
recorte de menciones por capitalización; no existe la taxonomía de cuatro
tipos, de modo que un enrutador léxico no tiene qué enrutar; y no hay
identificadores externos, de modo que la precisión de fusión de la memoria
canónica no puede medirse con oro.

\begin{table}[t]
\centering\footnotesize\setlength{\tabcolsep}{4pt}
\caption{Los dos bancos, tal como se usan aquí (subconjuntos de reproducción
de \hippo{}).}
\label{tab:diferencias}
\begin{tabular}{p{0.20\textwidth}p{0.36\textwidth}p{0.36\textwidth}}
\toprule
 & 2WikiMultihopQA & HotpotQA\\
\midrule
origen & 2020; plantillas sobre tripletes de Wikidata & 2018; preguntas escritas por anotadores\\
banco & $1\,000$ del \emph{dev}; corpus $6\,119$ & $1\,000$ del \emph{dev}, nivel \emph{hard}; corpus $9\,811$\\
tipos & comparación 244, composicional 413, inferencia 108, comparación-puente 235 & puente 811, comparación 189\\
pasajes oro & 2 (4 en comparación-puente) & 2\\
oro adicional & tripletes de evidencia con QIDs de Wikidata & frases de apoyo; sin identificadores\\
estilo & corto, regular, entidades en \emph{Title Case} & largo, descriptivo, minúsculas frecuentes\\
sondeo & 200 ($50$ por tipo) del \emph{dev} restante; mini-corpus $1\,558$ & 200 ($100$ por tipo) del \emph{train}, \emph{hard}; mini-corpus $1\,990$\\
\bottomrule
\end{tabular}
\end{table}

% =========================================================================
\section{Método: qué se conserva, qué se adapta}
\label{sec:adaptacion}

El sistema es el del documento de 2Wiki en su versión v4. Esta sección no lo
repite: recuerda cada etapa en una frase y detalla solo lo que se cambió
para que el traslado fuese posible sin ajuste. La regla que gobernó la
adaptación es de \emph{generalidad}: se descarta todo mecanismo que dependa
del idioma de las preguntas, de la forma de los títulos o de la taxonomía de
tipos de un conjunto concreto, aunque en ese conjunto rindiese.

\subsection{Indexación (sin cambios)}
\label{sec:indexacion}

\begin{enumerate}[leftmargin=1.6em, itemsep=1pt]
\item \textbf{Extracción reificada.} Una llamada a \texttt{gpt-4o-mini} por
  pasaje devuelve sus entidades y sus aserciones ---sujeto, relación, objeto,
  calificadores y una frase autocontenida sin pronombres---. Consigna
  idéntica a la de 2Wiki. En HotpotQA: $9\,811$ pasajes, $87\,629$
  aserciones ($8{,}9$ por pasaje) y $74\,947$ menciones de entidad, sin un
  solo error de formato.
\item \textbf{Memoria canónica.} Una ficha por (pasaje, mención) con su
  nombre y las frases en que participa; candidatos a la misma entidad por dos
  canales (coseno de nombre, coseno de ficha) con restricciones duras
  (fechas de nacimiento y muerte, ordinales); fusión automática por nombre
  idéntico o casi idéntico con contexto afín, juez de lenguaje para la banda
  ambigua, cierre transitivo. La versión v4 (sección~\ref{sec:memoria}) añade
  procedencia y curación. En HotpotQA: $77\,589$ fichas (las $74\,947$
  menciones más un título sintético por pasaje cuyo sujeto no fue extraído
  con su forma de título) $\to$ $48\,856$ entradas, $9\,525$ de ellas con más
  de una mención.
\item \textbf{Grafo canónico.} Nodos pasaje, entidad canónica y aserción;
  aristas aserción$\to$pasaje (procedencia), aserción$\to$entidad (rol de
  sujeto, de objeto o de participante $n$-ario) y entidad$\to$pasaje
  (mención). Sin aristas de sinonimia: las variantes ya colapsaron en la
  memoria. En HotpotQA: $146\,274$ nodos ($9\,811$ + $48\,834$ + $87\,629$) y
  $311\,979$ aristas. Con los roles invertidos para la consulta, el grafo es
  \emph{por capas} ---entidad $\to$ aserción $\to$ pasaje--- y, como se
  demostró en el documento de 2Wiki, el punto fijo del paseo tiene forma
  cerrada: el paseo no encadena saltos; los encadena la siembra.
\item \textbf{Almacenes vectoriales.} Una incrustación por aserción (su
  frase), por entrada (nombre, alias y descripción), por pasaje (título y
  texto) y, nueva en v4, por alias (el nombre solo, para el enlazador),
  todas con \texttt{text-embedding-3-small}.
\end{enumerate}

\subsection{Consulta: las tres reglas retiradas y sus sustitutos}
\label{sec:consulta}

La consulta de la versión evaluada en 2Wiki tenía tres dependencias del
conjunto. Se enumeran con lo que las sustituye.

\paragraph{(R1) Enrutador léxico $\to$ una sola estrategia.}
En 2Wiki, un clasificador por expresiones regulares asignaba cada pregunta a
uno de los cuatro tipos y decidía si la selección final de conjunto se
ejecutaba (no en comparación, donde la v3 ya saturaba). Nada de eso existe
en HotpotQA. En v4 hay una sola estrategia para toda pregunta: analista,
paseo, salto dirigido si hay huecos y selección de conjunto \emph{siempre}.
El clasificador sobrevive únicamente como etiqueta para el desglose de
2Wiki. La sección~\ref{sec:sondeo} mide el efecto de ejecutar el conjunto
también en comparación: $+2$ puntos de $\mathrm{FC@}5$ en ese tipo.

\paragraph{(R2) Recorte por mayúsculas $\to$ reconocimiento por el analista y enlazador de nombre a nombre.}
El enlazador de menciones convertía la pregunta en tramos por capitalización
y los casaba con los alias por igualdad tras plegado; un respaldo vectorial
contra la ficha ($\cos \ge 0{,}80$) estaba, medido, inactivo (3 de 275
tramos no exactos). En v4: (a) el \emph{analista} devuelve las menciones de
la pregunta con su forma canónica, y esas formas se enlazan con el mismo
enlazador ---es el reconocimiento de entidades que no depende de mayúsculas,
L1---; (b) el escalón 2 pasa a comparar \emph{nombre con nombre}: coeficiente
de Dice sobre trigramas de caracteres del alias plegado ($\ge 0{,}80$), que
mide truncamientos y erratas y penaliza la palabra que falta (\gl{}bad
education movie\gr{} no es \gl{}Bad Education\gr{}); el coseno de nombre
($\ge 0{,}85$) solo admite candidatos cuando el léxico no admite a nadie
(variantes sin superficie común); entre entradas a menos de $0{,}02$ de la
mejor decide la ficha con descripción (L2); (c) el escalón exacto prueba
primero la forma con el desambiguador entre paréntesis conservado (\gl{}X
(1922 film)\gr{} ya no colisiona con \gl{}X (1923 film)\gr{}) y las
colisiones de alias se resuelven por ficha y no por orden de carga (L4
parcial); (d) el posesivo se retira solo al final del tramo (\gl{}Lothair
II's mother\gr{}), porque dentro del tramo es parte del nombre (\gl{}God's
Gift to Women\gr{}). Un parche de prefijos (reintentar el alias anteponiendo
\gl{}the\gr{}) habría resuelto $168$ de los $275$ tramos fallidos de 2Wiki y
se descartó por no generalizable.

\paragraph{(R3) Candidatos del salto por subcadena $\to$ ancla explícita.}
El salto dirigido tomaba como candidatos estructurales las aserciones del
pasaje propio de la entidad enlazada cuyo nombre apareciese como subcadena
del texto del hueco. En v4 el analista devuelve, por cada hueco sin cubrir,
la \emph{entidad ancla} a la que el hueco se refiere (la película, la
persona), y los candidatos estructurales son las aserciones del pasaje propio
de esa ancla, enlazada por el diccionario; la subcadena queda como respaldo
cuando el ancla no enlaza.

\paragraph{Rondas de salto.}
Para conjuntos con cadenas de tres y cuatro saltos (MuSiQue), el salto
dirigido admite varias rondas: si tras un salto quedan huecos con comodín
$[X]$, el analista re-planifica con los hechos descubiertos a la vista y se
salta de nuevo, a lo sumo tres veces. En HotpotQA, de dos saltos, la
segunda ronda casi nunca se activa; se incluye porque la configuración se
congela para los tres conjuntos.

\paragraph{Demostraciones del analista.}
Las dos demostraciones de la consigna del analista se re-derivaron del
conjunto de \emph{calibración} de 2Wiki (\gl{}Convoy Buddies / Tip Toes\gr{},
\gl{}Zoolander 2\gr{}) en lugar de proceder del banco, lo que retira la
reserva de contaminación de diseño que el documento de 2Wiki mantenía sobre
el Plan A. Ninguna demostración procede de HotpotQA.

\subsection{Lo que no se ajustó}
\label{sec:congelado}

Para que la afirmación \gl{}sin ajuste\gr{} sea comprobable, la lista
completa de lo que HotpotQA heredó tal cual: los diales del paseo (peso de
la presa densa $0{,}05$, amortiguación $0{,}5$, ancla débil $0{,}05$, ancla
del enlazador $0{,}5$: los publicados por \hippo{}); los umbrales del
enlazador ($0{,}80$ Dice, $0{,}85$ coseno de nombre, $0{,}02$ de empate) y de
la memoria canónica ($0{,}75$ candidato, $0{,}92$ fusión automática de
ficha, $0{,}80$ nombre, $15$ vecinos); los tamaños de candidatos (Top-40,
frontera de $15$ por entidad con tope de $20$ nuevos, $16$ por hueco en el
salto, top-12 en el conjunto); las tres consignas y sus demostraciones; el
tope dinámico del analista ($\min(8,\max(2,|\text{plan}|))$); el modelo
(\texttt{gpt-4o-mini}) y el codificador. Lo único distinto de un conjunto a
otro son los datos.

\subsection{Banco de pruebas generalizado}
\label{sec:harness}

El harness normaliza cada conjunto a un esquema común (identificador,
pregunta, respuesta, tipo, títulos oro) y asigna un prefijo por conjunto a
los \emph{splits} y a los artefactos (\texttt{hotpot\_benchmark},
\texttt{hotpot\_sondeo}), de modo que extracción, memoria, grafo, recuperador
y medición aceptan cualquier conjunto sin cambios de código. MuSiQue queda
normalizado desde sus párrafos con marca de apoyo y su descomposición oro
(el tipo es el número de saltos), a falta de resolver que sus títulos se
repiten ($647$ de $11\,656$ pasajes), lo que exige definir el pasaje propio
como conjunto.

% =========================================================================
\section{Curación auditada de la memoria canónica (v4)}
\label{sec:memoria}

La memoria canónica se cura \emph{antes} de indexar el conjunto nuevo,
porque un índice nuevo heredaría los defectos de la consolidación. La
curación se valida en 2Wiki, el único conjunto con oro externo, y se aplica
sin cambios a HotpotQA.

\paragraph{Auditoría con oro externo.}
2Wiki trae en cada pregunta los identificadores Wikidata de las entidades de
sus evidencias. Tomando como unidad el par (forma, QID) de las entidades de
la cadena ---sujetos de evidencia y títulos oro; $2\,487$ pares, $2\,293$
resueltos a una entrada---, una entrada con dos QIDs de formas distintas es
una \emph{sobre-fusión} y un QID repartido en dos entradas una
\emph{infra-fusión}; una misma forma bajo dos QIDs se excluye como ruido del
oro. Registrando la \emph{procedencia} de cada unión (quién decidió que dos
fichas son la misma entidad) se atribuye cada sobre-fusión al origen que
tiende el puente entre los dos QIDs. En la versión v3 había $74$ entradas
sobre-fundidas de $2\,211$ tocadas ($3{,}3\%$) y $6$ infra-fusiones; el
origen puente más frecuente ($40$ de $74$) era la unión título$\to$sujeto
ejecutada cuando el título ya era mención propia del pasaje: unía al sujeto
con parientes, secuelas y homónimos.

\paragraph{Cuatro decisiones registradas.}
(i) La unión título$\to$sujeto solo es automática cuando el título es
mención \emph{sintética} y el sujeto candidato es único; si el título ya es
mención propia o hay empate, el par va al juez. (ii) La ficha-título sin
fechas hereda las del sujeto único del artículo, lo que reactiva la
restricción de fechas contra homónimos. (iii) Un nombre base idéntico o casi
idéntico con desambiguador entre paréntesis distinto va al juez en lugar de
fundirse. (iv) Cada entrada conserva el recuento de sus uniones por origen.
En 2Wiki la sobre-fusión baja de $74$ a $32$ entradas con la infra-fusión
intacta; el residuo son parientes con nombre casi igual que el propio juez
acepta. El coste en 2Wiki fue de unos céntimos ($2\,325$ adjudicaciones
nuevas); la curación cede dos preguntas del sondeo que la v3 acertaba por
fusiones erróneas afortunadas (un futbolista etiquetado como director por el
propio conjunto; un emperador fundido con su hijo) y se adopta igualmente,
porque es el principio correcto y porque es la que necesitan los conjuntos
sin identificadores.

\paragraph{En HotpotQA.}
La tabla~\ref{tab:indice} da las cifras de la consolidación. Sin QIDs no
hay auditoría: la precisión de fusión solo podría estimarse con una muestra
etiquetada a mano, que no se ha hecho (sección~\ref{sec:validez}). Lo que sí
se observa es que la proporción de decisiones que el juez rechaza es
parecida a la de 2Wiki (títulos-mención: $1\,975$ aceptadas de $2\,881$,
$69\%$; empates: $447$ de $472$).

\begin{table}[t]
\centering\footnotesize\setlength{\tabcolsep}{5pt}
\caption{Índice de HotpotQA (banco y sondeo). Juez: pares enviados $\to$
aceptados. Coste con \texttt{gpt-4o-mini} y \texttt{text-embedding-3-small}.}
\label{tab:indice}
\begin{tabular}{lrr}
\toprule
 & banco & sondeo\\
\midrule
pasajes & $9\,811$ & $1\,990$\\
aserciones reificadas & $87\,629$ & $17\,025$\\
menciones de entidad (fichas) & $77\,589$ & $14\,685$\\
fichas-título con fechas heredadas & $1\,041$ & $196$\\
fusiones automáticas (nombre idéntico / casi idéntico) & $102\,592$ / $658$ & $10\,685$ / $167$\\
título$\to$sujeto sintético (directas) & $2\,017$ & $387$\\
juez: vecinos & $22\,635 \to 6\,664$ & $2\,742 \to 811$\\
juez: título ya mención & $2\,881 \to 1\,975$ & $588 \to 391$\\
juez: empate de sujeto & $472 \to 447$ & $102 \to 98$\\
entradas canónicas & $48\,856$ & $10\,335$\\
grafo: nodos / aristas & $146\,274$ / $311\,979$ & $29\,345$ / $60\,176$\\
coste: extracción / juez / incrustaciones & $3{,}9$ / $0{,}4$ / $0{,}4$ USD & $0{,}8$ / $0{,}1$ / $0{,}1$ USD\\
\bottomrule
\end{tabular}
\end{table}

% =========================================================================
\section{Protocolo}
\label{sec:protocolo}

\begin{enumerate}[leftmargin=1.6em, itemsep=1pt]
\item \textbf{Hipótesis pre-registrada} (H1 del documento de 2Wiki): en
  HotpotQA, la versión v3 supera la cadena completa publicada de CatRAG
  ($80{,}4$) con los mismos modelos, porque el segmento de puente equivale
  al composicional de 2Wiki. Se contrasta con la v4, que la contiene.
\item \textbf{Configuración congelada} (sección~\ref{sec:congelado}) y
  \textbf{modelos de paridad}: \texttt{gpt-4o-mini} en extracción, juez,
  analista, salto y conjunto; \texttt{text-embedding-3-small} en toda
  incrustación. Son los modelos con los que CatRAG publica su tabla.
\item \textbf{Sondeo} de $200$ para comprobar que nada se rompe y para las
  ablaciones; \textbf{medición única} del banco; sin segunda pasada.
\item \textbf{Suelo} propio: BM25 sobre el mismo corpus (implementación
  independiente de la de CatRAG, por lo que las dos cifras de BM25 no son
  comparables entre sí; sí lo son con sus respectivos sistemas).
\item \textbf{Comparación}: pareada frente a BM25 (mismas preguntas); frente
  a \hippo{} y CatRAG, por las cifras publicadas con los mismos modelos, sin
  intervalo (no se reprodujo \hippo{} sobre este corpus).
\item \textbf{Trazabilidad}: cada llamada al modelo se guarda en una caché
  SQLite por (modelo, resumen del mensaje); toda medición es reproducible a
  coste cero y cada paso de la sesión queda en una bitácora con hora,
  comando, coste y resultado.
\end{enumerate}

% =========================================================================
\section{Resultados}
\label{sec:resultados}

\subsection{Sondeo y ablaciones}
\label{sec:sondeo}

La tabla~\ref{tab:sondeo} da el sondeo. El sistema completo alcanza
$\mathrm{R@}5$ $98{,}8$ y $\mathrm{FC@}5$ $97{,}5$ (puente $95{,}0$,
comparación $100$) frente a $77{,}2$ / $58{,}5$ de BM25; unas $3$ llamadas
por pregunta. Dos ablaciones reparten la contribución. Sin salto dirigido,
$\mathrm{FC@}5$ cae a $95{,}5$ (puente $91{,}0$): pareado, $+2{,}0$
$[+0{,}5,+4{,}0]$ para el sistema completo, que gana $4$ preguntas y no
pierde ninguna; también $\mathrm{R@}2$ ($+2{,}2$, significativo) y
$\mathrm{R@}5$ ($+1{,}0$, significativo). Sin selección de conjunto,
$\mathrm{R@}2$ cae de $87{,}0$ a $83{,}0$ ($+4{,}0$ $[+1{,}5,+6{,}8]$,
significativo: el conjunto ordena temprano), $\mathrm{FC@}5$ de $97{,}5$ a
$96{,}5$ (n.s.) y, en comparación, de $100$ a $98{,}0$: es la evidencia de
que ejecutar el conjunto también en comparación (R1) no era neutro.

\begin{table}[t]
\centering\footnotesize\setlength{\tabcolsep}{4pt}
\caption{Sondeo de HotpotQA ($200$; \texttt{gpt-4o-mini}). Pareado: diferencia
sistema completo $-$ ablación (IC del $95\%$; gana / pierde).}
\label{tab:sondeo}
\begin{tabular}{lcccccl}
\toprule
sistema & $\mathrm{R@}2$ & $\mathrm{R@}5$ & $\mathrm{FC@}2$ & $\mathrm{FC@}5$ & $\mathrm{FC@}10$ & $\mathrm{FC@}5$ puente / comparación\\
\midrule
BM25 (suelo) & $58{,}0$ & $77{,}2$ & $27{,}5$ & $58{,}5$ & $83{,}0$ & $48{,}0$ / $69{,}0$\\
\sys{} v4 completo & $\mathbf{87{,}0}$ & $\mathbf{98{,}8}$ & $\mathbf{75{,}5}$ & $\mathbf{97{,}5}$ & $98{,}0$ & $95{,}0$ / $100$\\
\quad sin salto dirigido & $84{,}8$ & $97{,}8$ & $71{,}5$ & $95{,}5$ & $97{,}0$ & $91{,}0$ / $100$\\
\quad sin selección de conjunto & $83{,}0$ & $98{,}2$ & $68{,}0$ & $96{,}5$ & $98{,}5$ & $95{,}0$ / $98{,}0$\\
\midrule
\multicolumn{7}{l}{pareado completo $-$ sin salto: $\mathrm{FC@}5$ $+2{,}0$ $[+0{,}5,+4{,}0]$ (4/0); $\mathrm{R@}2$ $+2{,}2$ $[+0{,}2,+4{,}5]$; $\mathrm{R@}5$ $+1{,}0$ $[+0{,}2,+2{,}0]$}\\
\multicolumn{7}{l}{pareado completo $-$ sin conjunto: $\mathrm{R@}2$ $+4{,}0$ $[+1{,}5,+6{,}8]$ (17/4); $\mathrm{FC@}5$ $+1{,}0$ $[-1{,}0,+3{,}0]$ n.s.\ (3/1)}\\
\bottomrule
\end{tabular}
\end{table}

\subsection{Banco}
\label{sec:banco}

La tabla~\ref{tab:banco} sitúa la medición única del banco junto a las
referencias publicadas con los mismos modelos. \sys{} v4 obtiene
$\mathrm{R@}2$ $77{,}1$, $\mathrm{R@}5$ $94{,}7$ $[93{,}6,95{,}8]$,
$\mathrm{FC@}5$ $90{,}2$ $[88{,}2,92{,}0]$ y $\mathrm{FC@}10$ $93{,}6$; por
tipo, puente $93{,}6$ / $88{,}2$ y comparación $99{,}5$ / $98{,}9$
($\mathrm{R@}5$ / $\mathrm{FC@}5$). Frente a CatRAG, $+5{,}2$ de
$\mathrm{R@}5$ y $+9{,}8$ de cadena completa; frente a \hippo{}, $+7{,}6$ y
$+14{,}7$; frente al mejor codificador denso del mismo tamaño, $+13{,}4$ y
$+25{,}3$. Solo \hippo{} con NV-Embed-v2 (siete mil millones de parámetros,
GPU) publica un $\mathrm{R@}5$ mayor ($96{,}3$), sin cadena completa
informada. Pareado frente a nuestro BM25 sobre el mismo corpus: $\mathrm{FC@}5$
$+40{,}7$ $[+37{,}1,+44{,}1]$ (gana $432$, pierde $25$), puente $+42{,}8$,
comparación $+31{,}7$; $\mathrm{R@}5$ $+21{,}9$ $[+20{,}0,+23{,}8]$. La
hipótesis H1 se confirma con margen: pedía superar $80{,}4$ y la v4 lo
supera en diez puntos. El coste de la medición fue de $2\,982$ llamadas
($\approx 3$ por pregunta, $\approx 1$ USD, $62$ minutos).

\begin{table}[t]
\centering\footnotesize\setlength{\tabcolsep}{5pt}
\caption{Banco de HotpotQA ($1\,000$ preguntas, corpus $9\,811$). Referencias
de~\cite{catrag} (su tabla con \texttt{gpt-4o-mini} +
\texttt{text-embedding-3-small}; FCR $=\mathrm{FC@}5$; F1 con lector
Llama-3.3-70B, no medido aquí). Última fila de referencia: \hippo{} con
NV-Embed-v2 de~\cite{hipporag2}.}
\label{tab:banco}
\begin{tabular}{lcccl}
\toprule
sistema & $\mathrm{R@}2$ & $\mathrm{R@}5$ & $\mathrm{FC@}5$ & nota\\
\midrule
BM25 (CatRAG) & --- & $64{,}8$ & $38{,}3$ & F1 $54{,}1$\\
Contriever~\cite{contriever} & --- & $75{,}3$ & $54{,}1$ & \\
GTR~\cite{gtr} & --- & $73{,}9$ & $53{,}0$ & \\
denso \texttt{text-embedding-3-small} & --- & $81{,}3$ & $64{,}9$ & \\
RAPTOR~\cite{raptor} & --- & $79{,}5$ & $61{,}4$ & \\
\hippo{}~\cite{hipporag2} & --- & $87{,}1$ & $75{,}5$ & F1 $69{,}4$\\
CatRAG~\cite{catrag} & --- & $89{,}5$ & $80{,}4$ & F1 $71{,}4$\\
\midrule
BM25 (nuestra implementación, mismo corpus) & $56{,}1$ & $72{,}8$ & $49{,}5$ & suelo del pareado\\
\textbf{\sys{} v4 (este trabajo)} & $\mathbf{77{,}1}$ & $\mathbf{94{,}7}$ & $\mathbf{90{,}2}$ & $\mathrm{FC@}10$ $93{,}6$; puente $88{,}2$, comparación $98{,}9$\\
\midrule
\hippo{} con NV-Embed-v2 (7B) & --- & $96{,}3$ & --- & techo de hardware\\
\bottomrule
\end{tabular}
\end{table}

\subsection{Del sondeo al banco, y de 2Wiki a HotpotQA}
\label{sec:comparacion}

La caída del sondeo al banco ($\mathrm{FC@}5$ $97{,}5 \to 90{,}2$) es la del
corpus cinco veces mayor: más distractores con títulos parecidos, más
homónimos, más hechos afines por coseno. Con la misma configuración y el
mismo modelo, 2Wiki da $98{,}1$ / $96{,}0$ y HotpotQA $94{,}7$ / $90{,}2$. La
diferencia es coherente con la naturaleza de los conjuntos: las preguntas
de 2Wiki nombran sus entidades con la forma del título y encadenan
relaciones de una lista corta (director, madre, fecha), de modo que el
enlazador ancla casi siempre y el analista reconoce el patrón; las de
HotpotQA describen (\gl{}the radio station that serves the area comprising
all of Belknap County\gr{}), nombran en minúsculas y exigen a veces un puente
que el propio pasaje oro solo insinúa. Aun así, el segmento de puente de
HotpotQA ($88{,}2$) queda en la banda del composicional de 2Wiki ($95{,}9$)
más que en la de \hippo{} en ese mismo conjunto ($75{,}5$), que era lo que H1
predecía.

% =========================================================================
\section{Análisis de fallos}
\label{sec:fallos}

De las $1\,000$ preguntas, $98$ no consiguen la cadena completa a $k{=}5$:
$96$ de puente y $2$ de comparación. En $90$ de ellas falta \emph{un solo}
pasaje ---el puente---, que queda en los puestos 6--10 en $37$ casos, en
11--20 en $23$ y fuera del top-20 en $46$; en $8$ faltan los dos. Las $35$
preguntas sin ninguna mayúscula interior fallan algo más ($6$, un $17\%$,
frente al $9{,}8\%$ global): el recorte por capitalización, que en v4 solo
alimenta la frontera del analista, todavía pesa, y la vía es que el
reconocimiento de entidades del analista preceda también a la frontera (L1
completo).

Los fallos leídos uno a uno en el sondeo no muestran un patrón nuevo:
oro secundario que el pasaje principal apenas conecta (\gl{}V8 engine\gr{} $\to$
\gl{}Chevrolet Camaro (first generation)\gr{}, en el puesto 19), homónimos de
serie (\gl{}ABU TV Song Festival 2014\gr{} entre las ediciones de 2012, 2013 y
2016), puentes descritos sin nombre (\gl{}the man who called Log College's
structure very plain in his journal\gr{} $\to$ George Whitefield, ausente) o
un mascota cuya especie no aparece en el pasaje de la universidad
(\gl{}Scrappy Moc\gr{} $\to$ \gl{}Northern mockingbird\gr{}). Son los mismos
tres mecanismos de fallo que el documento de 2Wiki cataloga ---el filtro no
selecciona el puente aunque lo ve, el enlazador no ancla, el hecho puente no
está verbalizado en el pasaje--- con la proporción desplazada hacia el
tercero, propio de un conjunto donde el puente lo escribió una persona y no
una plantilla.

% =========================================================================
\section{Discusión}
\label{sec:discusion}

\paragraph{Por qué generaliza.}
Tres decisiones de diseño explican que el sistema no necesitara ajuste. La
\emph{memoria canónica} convierte la variación de superficie ---el problema
que en HotpotQA se agrava con las menciones descriptivas--- en un problema
de indexación resuelto una vez, y da a cada entidad un pasaje propio al que
la siembra puede depositar masa directamente. El \emph{analista} sustituye
al enrutador y al filtro: no clasifica la pregunta, la \emph{planifica}, y
un plan de huecos con anclas es el mismo objeto en cualquier conjunto. Y el
\emph{salto dirigido} es la única pieza que encadena de verdad, porque el
grafo por capas no lo hace; su contribución, significativa en los dos
sondeos, es la que separa a este sistema de los paseos con reinicio sobre
grafos de entidades.

\paragraph{Qué costó.}
Indexar HotpotQA costó $\approx 4{,}7$ USD (extracción $3{,}9$, juez $0{,}4$,
incrustaciones $0{,}4$) y unas dos horas de reloj con ocho hilos; consultar,
$\approx 0{,}001$ USD por pregunta con \texttt{gpt-4o-mini}. La memoria
canónica es la pieza más cara en tiempo (producto de $77\,589 \times 77\,589$
cosenos por bloques y $3\,240$ llamadas al juez), y la que más se
beneficiaría de los candidatos por umbral sobre grupos de nombre plegado
que el documento de 2Wiki propone (E5).

\paragraph{Qué se aprendió de la curación.}
La auditoría con oro externo enseñó que una parte de las fusiones erróneas
\emph{ayudaba} en 2Wiki: fundir a un emperador con su hijo hace alcanzable
el pasaje del hijo desde el del padre, que es justo lo que pide una pregunta
de parentesco. Es una ganancia por accidente, no por diseño, y no se puede
esperar de un conjunto sin esa estructura familiar; renunciar a ella para
tener una memoria precisa y auditable es la decisión correcta, y HotpotQA
es el primer conjunto que la valida.

% =========================================================================
\section{Consideraciones de validez y limitaciones}
\label{sec:validez}

\begin{enumerate}[leftmargin=1.6em, itemsep=1pt]
\item \textbf{Sin reproducción propia de \hippo{} ni de CatRAG sobre este
  corpus.} La comparación con ellos es por cifras publicadas con los mismos
  modelos ---en 2Wiki, nuestra reproducción de \hippo{} coincidió al decimal
  con la del artículo de CatRAG, lo que da confianza en el protocolo---,
  pero sin intervalo pareado. Reproducir \hippo{} sobre HotpotQA cuesta
  unos $5$ USD y queda pendiente.
\item \textbf{Sin auditoría de la memoria canónica.} HotpotQA no trae
  identificadores externos; la precisión de fusión de las $48\,856$ entradas
  no está medida. La curación se validó en 2Wiki; su efecto aquí se infiere.
\item \textbf{Un solo modelo y una sola medición.} La fila de paridad
  (\texttt{gpt-4o-mini}) es la única del banco; en 2Wiki, el modelo adoptado
  en producción (\texttt{deepseek-chat}) rinde un punto más y su medición en
  HotpotQA está pendiente. La medición única evita el ajuste al banco pero
  no informa de la varianza entre ejecuciones (en 2Wiki, dos ejecuciones
  del mismo sistema difirieron en dos preguntas de mil).
\item \textbf{Oro a nivel de título.} Las frases de apoyo no se usan; un
  pasaje oro recuperado por una frase que no es la de apoyo cuenta como
  acierto. Es el protocolo de \hippo{} y CatRAG, y por eso se conserva.
\item \textbf{Solo nivel \emph{hard}.} El subconjunto de reproducción es
  íntegramente \emph{hard}; nada se dice de los niveles \emph{easy} y
  \emph{medium}.
\item \textbf{Sin respuesta.} Se mide recuperación, no la respuesta final
  (F1, tasa de éxito conjunto); la columna F1 de la tabla~\ref{tab:banco}
  queda vacía hasta que se añada el lector.
\item \textbf{Independencia de las decisiones.} Todas las decisiones de v4 se
  tomaron sobre el sondeo de 2Wiki y su conjunto de calibración; ninguna
  sobre datos de HotpotQA. Las ablaciones de la sección~\ref{sec:sondeo} se
  midieron \emph{después} de fijar la configuración y no la modificaron.
\end{enumerate}

% =========================================================================
\section{Trabajo futuro}
\label{sec:futuro}

En orden de valor: (1) la medición de los dos bancos con el modelo adoptado
(\texttt{deepseek-chat}) y la reproducción de \hippo{} sobre HotpotQA para
completar los pareados; (2) MuSiQue, que trae cadenas de tres y cuatro
saltos con su descomposición oro y exige antes el pasaje propio como
conjunto de pasajes con el mismo título; allí se contrastan H2 (la ventaja
se concentra en tres y cuatro saltos y depende de las rondas de salto) y H3
(la calidad del plan del analista, medida contra la descomposición oro,
predice la cadena completa mejor que la afinidad de la pregunta); (3) el
lector y las métricas de respuesta, para comparar con los F1 y JSR
publicados; (4) L1 completo ---reconocimiento de entidades antes de la
frontera--- motivado por las preguntas en minúsculas; (5) los candidatos por
umbral en la memoria canónica, para abaratar la indexación de corpus
mayores.

% =========================================================================
\section*{Reproducibilidad}

Código en la rama \texttt{eval-wiki2} del repositorio (módulos
\texttt{eval/wiki2.py}, \texttt{ingest/wiki2\_openie.py},
\texttt{ingest/wiki2\_entidades.py}, \texttt{ingest/wiki2\_grafo.py},
\texttt{retrieval/wiki2\_catrag.py}, \texttt{retrieval/wiki2\_plan.py},
\texttt{eval/wiki2\_correr.py}); artefactos con sufijo \texttt{\_v4} en
\texttt{artifacts/wiki2/}; resultados por pregunta en
\texttt{resultados\_plan\_v4\_4omini\_hotpot\_benchmark.jsonl} y agregados y
pareados junto a ellos; bitácora de la sesión en
\texttt{BITACORA\_2026-08-30.md}; cuaderno de trazas en
\texttt{notebooks/recuperacion\_hotpotqa.ipynb}. Comandos:
\texttt{wiki2 --generar-splits --dataset hotpot --n-por-tipo 100};
\texttt{wiki2\_openie --split hotpot\_benchmark};
\texttt{wiki2\_entidades --split hotpot\_benchmark --ver \_v4};
\texttt{wiki2\_grafo --split hotpot\_benchmark --canonico --ver \_v4};
\texttt{wiki2\_correr --split hotpot\_benchmark --modelo gpt-4o-mini --pareado bm25\_hotpot\_benchmark}.

% =========================================================================
\begin{thebibliography}{99}\small

\bibitem{hotpotqa} Z.~Yang, P.~Qi, S.~Zhang, Y.~Bengio, W.~W.~Cohen,
R.~Salakhutdinov, C.\,D.~Manning. HotpotQA: A dataset for diverse,
explainable multi-hop question answering. \emph{EMNLP}, 2018.

\bibitem{wiki2} X.~Ho, A.-K.~Duong Nguyen, S.~Sugawara, A.~Aizawa.
Constructing a multi-hop QA dataset for comprehensive evaluation of
reasoning steps. \emph{COLING}, 2020.

\bibitem{musique} H.~Trivedi, N.~Balasubramanian, T.~Khot, A.~Sabharwal.
MuSiQue: Multihop questions via single-hop question composition.
\emph{TACL}, 2022.

\bibitem{hipporag} B.~Jiménez Gutiérrez, Y.~Shu, Y.~Gu, M.~Yasunaga,
Y.~Su. HippoRAG: Neurobiologically inspired long-term memory for large
language models. \emph{NeurIPS}, 2024.

\bibitem{hipporag2} B.~Jiménez Gutiérrez, Y.~Shu, W.~Qi, S.~Zhou, Y.~Su.
From RAG to memory: Non-parametric continual learning for large language
models. 2025. arXiv:2502.14802.

\bibitem{catrag} K.\,H.~Lau, F.~Zhang, B.~Ruan, Y.~Zhou, Q.~Guo,
R.~Zhang, X.~Zhou. Breaking the static graph: Context-aware traversal for
graph-based RAG. \emph{Findings of ACL}, págs.~5849--5863, 2026.

\bibitem{raptor} P.~Sarthi, S.~Abdullah, A.~Tuli, S.~Khanna, A.~Goldie,
C.\,D.~Manning. RAPTOR: Recursive abstractive processing for
tree-organized retrieval. \emph{ICLR}, 2024.

\bibitem{contriever} G.~Izacard, M.~Caron, L.~Hosseini, S.~Riedel,
P.~Bojanowski, A.~Joulin, E.~Grave. Unsupervised dense information
retrieval with contrastive learning. \emph{TMLR}, 2022.

\bibitem{gtr} J.~Ni, C.~Qu, J.~Lu, Z.~Dai, G.~Hernández Ábrego, J.~Ma,
V.~Zhao, Y.~Luan, K.~Hall, M.-W.~Chang, Y.~Yang. Large dual encoders are
generalizable retrievers. \emph{EMNLP}, 2022.

\bibitem{nvembed} C.~Lee, R.~Roy, M.~Xu, J.~Raiman, M.~Shoeybi,
B.~Catanzaro, W.~Ping. NV-Embed: Improved techniques for training LLMs as
generalist embedding models. \emph{ICLR}, 2025.

\bibitem{bm25} S.~Robertson, H.~Zaragoza. The probabilistic relevance
framework: BM25 and beyond. \emph{Foundations and Trends in Information
Retrieval}, 3(4):333--389, 2009.

\bibitem{pprsurvey} M.~Yang, H.~Wang, Z.~Wei, S.~Wang, J.-R.~Wen.
Efficient algorithms for personalized PageRank computation: A survey.
\emph{IEEE TKDE}, 2024.

\end{thebibliography}

\end{document}
